\scanpage{398}\mspage{122}
localiseurs fondamentaux généraux !). En
effet, $W$ se récupère de façon unique
à l'aide de $D=\mathrm{HOT}_{W}$, ou de \unsure{tout}
dérivateur équivalent à celui-ci,
comme l'ens des $D$-équivalences dans
$\mathrm{Cat}$:$^{\ast}$ c'est \unsure{la|le} localiseur fondamental
% slanted
\margin{Mais ceci n'est prouvé que moyennant $\mathrm{Loc}(3\ \mathrm{bis})$}
(correspondant à une notion d'équivalence
``universelle'', i.e. dans la catégorie $\mathrm{Cat}$
de tous les ``\unsure{engins|espaces}'' \deleted{\unsure{abs}}\unsure{élémentaires})
canoniquement associé à \unsure{ce|un} dérivateur.
Ce n'est donc nullement un localiseur
de référence \unsure{plus ou moins fin|fixé une fois}, tel $\mathbb{W}_{\infty}$,
qui est apte à jouer ce rôle !

Cela nous amène à nouveau à une
question du genre déjà soulevé : Si
$W$ \inserted{$\subset\mathbb{W}_{\infty}$} est le plus petit localiseur qui
satisfait à $\mathrm{Loc}(6'\mathrm{a},8)$ \inserted{(ou $\mathrm{Loc}(3\ \mathrm{bis},6'\mathrm{a},8)$)}, peut-on
dire a) que $W$ soit $=$ \uline{spécial} i.e.
% archive sheet number: 397
